\chapter{Conclusion générale}

Au moment de conclure, ce que je retiens de ces deux années chez Souchier-Boullet ne tient pas dans la liste des livrables. En travaillant à la fois sur la maintenance des modèles Autodesk Inventor et sur le développement de configurateurs SQL, j'ai vu de l'intérieur comment se construit, se formalise et se sécurise une chaîne numérique industrielle, avec ses compromis, ses fragilités et ses utilisateurs qui n'attendent pas que l'outil soit parfait pour en avoir besoin.

Le travail mené montre qu'une transformation numérique pertinente ne passe pas uniquement par l'ajout d'outils nouveaux. Elle repose surtout sur la capacité à rendre explicites des règles métier, à structurer des données jusque-là dispersées, à documenter les interfaces entre systèmes et à concevoir des solutions qui restent lisibles pour les équipes. C'est dans cette logique que le configurateur développé a pris sa valeur : non seulement comme outil de configuration, mais comme support de capitalisation et de dialogue interservices.

Le mémoire met également en évidence une idée qui me paraît essentielle pour la suite : dès que l'on cherche à relier plusieurs systèmes, la question centrale n'est pas seulement \og comment faire communiquer les outils \fg{}, mais \og comment s'assurer qu'ils se comprennent encore lorsque l'un d'eux évolue \fg{}. Le chapitre scientifique consacré au contrat d'interface formalise précisément ce point et propose une piste concrète pour renforcer la robustesse de la liaison entre configurateur SQL et CAO paramétrique.

Sur le plan personnel, cette expérience a confirmé mon intérêt pour les sujets situés à l'interface entre ingénierie produit, structuration de données et outils numériques. Elle m'a aussi appris qu'une solution réussie est rarement celle qui impressionne le plus sur le plan technique, mais plutôt celle qui rend le travail collectif plus fiable, plus clair et plus durable. C'est dans cette direction que je souhaite désormais poursuivre mon parcours.
